% ============================================================
%  path_summary.tex
%  Short overview of the route from the manuscript's unresolved
%  Experiment 2 to the current state.  Deliberately compact:
%  one page of narrative, the rest tables.
% ============================================================
\documentclass[11pt,a4paper]{article}

\usepackage{amsmath,amssymb}
\usepackage{geometry}
\geometry{margin=2.3cm}
\usepackage{booktabs}
\usepackage{array}
\usepackage{xcolor}
\usepackage{microtype}
\usepackage{parskip}

\title{From a failed experiment to a working method:\\the route taken}
\author{}
\date{}

\begin{document}
\maketitle
\vspace{-1.5em}

\section*{1. Where the manuscript stops}

Table~3 of the manuscript reports Experiment~2 with DCL at $4.4$--$9.3\times$
the RVI optimum, and the text notes: \emph{``To be fixed in the future.''}
That is the starting point.  The finding of the work since is that this was not
a tuning problem: it was five distinct mismatches between the reinforcement
learning abstraction and the semi-Markov structure of the model, four of which
turned out to share a single root cause.

\section*{2. Five failure mechanisms, each isolated by ablation}

Every fix was validated by removing it from the otherwise complete recipe and
re-running 8 seeds on Experiment~2 (binary cost).  ``Within 5\%'' is measured
against the exact RVI optimum.

\begin{center}
\begin{tabular}{p{4.0cm}p{6.6cm}cc}
\toprule
\textbf{Mechanism} & \textbf{What goes wrong} & \textbf{Without} & \textbf{Failure mode} \\
\midrule
Semi-MDP discounting & Decisions are separated by a variable number of periods,
including zero.  Discounting each decision as if a period passed subsidises
idling, most strongly where the candidate list is longest. & 0/8 & never-serve \\[2pt]
Value-target scale & Returns are $O(10^2$--$10^3)$, advantages $O(1)$; the value
loss monopolises the shared trunk. & 0/8 & exactly bimodal:
never-serve or FIFO clone \\[2pt]
Data distribution & One deep-backlog excursion moves all environments into a
region where the shaping term saturates; behaviour on deployment-relevant
states degrades unobserved. & 1/8 & never-serve \\[2pt]
Reward sparsity & Binary cost is flat before the deadline; the penalty arrives
long after the decision that caused it. & 2/8 & never-serve \\[2pt]
Extraction & The deployed policy is the \emph{argmax}; training only certifies
the \emph{stochastic} policy. & 3/8 & argmax $\gg$ stochastic \\
\midrule
\textbf{Full recipe} & & \textbf{6/8} & \\
\bottomrule
\end{tabular}
\end{center}

The two attractors recur exactly across every failed cell --- never-serve at
$9.876\times$ and the FIFO clone at $1.738\times$ the optimum --- which is what
makes these diagnoses rather than symptoms.

\section*{3. The root cause: the action decomposition}

Mechanism~5 is the interesting one.  In the sequential candidate-queue
formulation a skipped job returns at the next tick.  If a state offers $m$
further presentations and the policy serves with probability $p$, the job is
served in time with probability $1-(1-p)^{m+1}$ --- close to one even for $p$ well
below $\tfrac12$.  The stochastic policy is therefore behaviourally fine where
the logit sign is wrong, the advantage vanishes, and training never corrects
it; but the argmax turns exactly those states into deterministic errors.  Most
of the stabilisation apparatus (temperature annealing, snapshots, readout
selection, seed selection) existed only to contain this.

Reformulating so that each idle capacity unit makes one $(N{+}1)$-way decision
(idle, or serve type $a$) removes it at the source.  The two formulations admit
the same physical policies --- FIFO and never-serve reproduce bit-exactly, and
the RVI optimum agrees on both instances --- so only the credit structure
changes.  This is also what the supervisor's own comment on \S3.2 proposed.

\begin{center}
\begin{tabular}{lcc}
\toprule
Experiment 2, binary cost & sequential space & per-event space \\
\midrule
with temperature annealing & 6/8 & 4/8 \\
\textbf{without} temperature annealing & 3/8 & \textbf{8/8, exact optimum} \\
\bottomrule
\end{tabular}
\end{center}

The apparatus is not merely unnecessary in the new space --- it is harmful there.
On Experiment~3 the dissociation runs the other way (7/8 with annealing, 3/8
without), so the honest statement is that a stabiliser is matched to a failure
mode and outlives its usefulness with it.

\section*{4. Three objectives}

\begin{center}
\begin{tabular}{p{3.0cm}p{5.4cm}p{3.4cm}cc}
\toprule
\textbf{Objective} & \textbf{Definition} & \textbf{Needs} & \textbf{Exp 2} & \textbf{Exp 3} \\
\midrule
Binary / FIL lateness & $c\,\mathbf{1}\{W>D\}$ per tick & shaping; small batch & 8/8 & 7/8 \\
Queue lateness & $c\sum_j (A_j-D)^+$ per tick, with conditional tail for
untracked jobs & average-reward; $64\times$ batch & 8/8 & 8/8 \\
Fraction late (tardiness flux) & $c$ once, at the deadline crossing & escalation
rule & 8/8 & 8/8 \\
\bottomrule
\end{tabular}
\end{center}

Seeds within 5\% of RVI, per-event action space, no per-cell tuning.  The
recipe follows from the tail of the cost function: queue lateness is the only
objective with unbounded per-tick cost ($\Theta(B^2)$ in the excess $B$), and
correspondingly the only one requiring the large-batch average-reward
configuration.

Two by-products worth recording.  \emph{Reward shaping} is potential-based and
therefore cost-neutral: FIFO and RVI evaluate identically with and without it to
four decimals, so the benchmark is preserved.  \emph{The escalation rule} ---
overdue work is served with forced priority --- is what makes the fraction-late
objective well posed: without it, letting a job rot suppresses its own arrival
stream in the FIL projection, which RVI slowly discovers, so learned policies
``beat'' the benchmark.  Escalation is implemented purely by restricting the
admissible action set to a singleton, so no algorithm was modified.

\section*{5. What the benchmark itself required}

RVI is the yardstick for all of the above, and it needed correcting twice.  The
truncation growth criterion ($1\%$ relative change in $g^*$) is adequate for the
binary cost, which saturates in the waiting level, but stops too early for
tail-heavy objectives: tightening it to $10^{-3}$ moved the queue-lateness
reference on Experiment~2 from $2492$ to $2350$.  Both corrections were
triggered by the same diagnostic --- a learned policy evaluating \emph{below} the
computed optimum, which can only mean a loose benchmark.

\section*{6. Current state and open questions}

Scaling is the active question, and the most recent finding is that it is a
modelling problem before it is an algorithmic one.  On a six-type chain the
learned policy beat the best heuristic by only $2.9\%$; the reason is not
under-training (a $2.5\times$ budget gives the same answer) but that on that
instance the myopic $c\mu$ rule is nearly optimal, because cost, service rate
and deadline all increase together, so the $c\mu$ ordering coincides with the
urgency ordering.

A screening of four candidate structures at three job types, where RVI is still
computable, shows what actually creates headroom:

\begin{center}
\begin{tabular}{lccc}
\toprule
Structure (3 types, load 2.0) & FIFO/RVI & c-$\mu$/RVI & $\min$ \\
\midrule
\textbf{cost and deadline anti-aligned} & \textbf{1.49} & \textbf{1.77} & \textbf{1.49} \\
equal costs, deadlines heterogeneous & 1.10 & 1.16 & 1.10 \\
non-interval skill overlap & 1.10 & 1.36 & 1.10 \\
generalist bottleneck (rewards idling) & 1.21 & 1.01 & 1.01 \\
\midrule
existing chain (all quantities aligned) & 1.19 & 1.19 & 1.19 \\
\bottomrule
\end{tabular}
\end{center}

Headroom comes from \emph{anti-alignment between cost and deadline}, not from
instance size, topological complexity, or idling opportunity.

Scaling that structure to four and six job types, with the recipe unchanged and
8 seeds per instance, gives the following.

\begin{center}
\begin{tabular}{lcccccc}
\toprule
Instance & types & FIFO & c-$\mu$ & PPO (8 seeds) & vs FIFO & vs c-$\mu$ \\
\midrule
anti-aligned, 3 & 3 & $1.372$ & $1.626$ & $0.929$--$0.965$ & $-30\%$ & $-43\%$ \\
anti-aligned, 4 & 4 & $1.518$ & $1.837$ & $1.124$--$1.201$ & $-21$ to $-26\%$ & $-35$ to $-39\%$ \\
anti-aligned, 6 & 6 & $1.745$ & $2.124$ & $1.422$--$1.634$ & $-6$ to $-18\%$  & $-23$ to $-31\%$ \\
\bottomrule
\end{tabular}
\end{center}

On the three-type instance, where the exact optimum is computable
($\mathrm{RVI} = 0.920$), all 8 seeds land within $1.0$--$4.9\%$ of it, with no
catastrophic seed.  The decisive comparison is at equal size: on the original
six-type chain the learned policy improved on the best heuristic by $2.9\%$,
whereas on the redesigned six-type instance it improves by $23$--$31\%$.  The
algorithm and the recipe are identical; only the alignment of cost, service
rate and deadline differs.  The earlier ceiling was a property of the
benchmark, not of the method.

Two honest qualifications.  The margin still narrows with instance size
($-43\%$, $-38\%$, $-31\%$ against c-$\mu$), and the seed spread widens (from
$0.9\%$ at three types to $15\%$ at six).  Moreover the two worst seeds across
the family are exactly the two for which the argmax evaluates worse than the
stochastic policy --- the extraction signature that the per-event reformulation
largely removes, but which leaves a trace at larger sizes.

\paragraph{Open.}
(i) A deadline-aware baseline (EDF / least slack, and a generalised $c\mu$
index) is still missing --- FIFO and $c\mu$ are both deadline-blind, which makes
them a weak comparison for a deadline objective.
(ii) Beyond three job types there is no exact reference; a fluid or LP lower
bound would be needed to say how much is left on the table.
(iii) The network is a fixed-width MLP, which assigns a separate semantic role
to each input position and so does not transfer across instance sizes;
parameter sharing over job types and server pools is the natural next step.

\end{document}
